% ===========================================================================
%  Review of the Obsidian notes (companion to expose_draft.tex, not submitted)
%  1  Errors and inaccuracies in the notes
%  2  Ratings of the terminology notes
%  3  Bibliography remarks
%  4  Answers to "Next Meeting.md"
%  5  Claude's ideas that do not belong in the exposé
% ===========================================================================

\documentclass[11pt,a4paper]{article}

\usepackage[T1]{fontenc}
\usepackage{lmodern}
\usepackage[english]{babel}
\usepackage{microtype}
\usepackage{csquotes}
\usepackage[a4paper,margin=2.3cm]{geometry}
\usepackage{parskip}
\usepackage{titlesec}
\titlespacing*{\section}{0pt}{1.2ex plus .5ex}{0.5ex}
\titlespacing*{\subsection}{0pt}{0.8ex plus .3ex}{0.3ex}
\usepackage{amsmath}
\usepackage{booktabs,tabularx}
\usepackage{enumitem}
\setlist{nosep,leftmargin=*}
\usepackage{xcolor}

\usepackage[backend=biber,style=numeric-comp,sorting=none,
            maxbibnames=3,giveninits=true,doi=false]{biblatex}
\addbibresource{references.bib}
\renewcommand*{\bibfont}{\small}
\AtEveryBibitem{%
  \iffieldundef{eprint}{}{\clearfield{url}}%
  \clearfield{urlyear}\clearfield{urlmonth}\clearfield{urlday}}

\usepackage[hidelinks]{hyperref}

\newcommand{\note}[1]{\textsf{#1}}                 % name of an Obsidian note
\newcommand{\sev}[1]{\textbf{[#1]}}                 % severity tag
\newcommand{\claim}[1]{\enquote{\textit{#1}}}       % quoted claim from a note
\newcommand{\fixed}{\textbf{\textcolor{green!50!black}{[corrected in the note]}}} % done

\begin{document}

\begin{center}
  {\Large\bfseries Review of the Obsidian Notes}\\[0.6ex]
  Companion to the exposé draft \textbullet\ \today
\end{center}

This document collects what did not belong in \texttt{expose\_draft.tex}: claims in
the notes that are wrong or inaccurate (Section~\ref{sec:errors}), a rating
of the notes that only explain terminology (Section~\ref{sec:ratings}),
remarks on the bibliography (Section~\ref{sec:bib}), and short answers to the
open questions in \note{Next Meeting.md} (Section~\ref{sec:meeting}).
Every claim was checked against the paper itself (arXiv abstract or full
text), not against secondary summaries.

Severity: \sev{Error} the claim is wrong; \sev{Inaccurate} partly right, but
misleading if copied into the exposé; \sev{Minor} wording, links, spelling.

\textbf{Status (11 Nov 2026):} every item in Section~\ref{sec:errors} has
been corrected directly in the notes (\fixed). None of these errors appear
in \texttt{expose\_draft.tex}. Three of the notes now live in the vault folder
\note{unimportant or old} (Integrated Gradients, Benchmarks for Mechanistic
Interpretability, and the Hubinger note); they were corrected there.

% ===========================================================================
\section{Errors and Inaccuracies in the Notes}\label{sec:errors}

\subsection{Notes linked from Structure.md}

\paragraph{\note{Origins - where geometry comes from.md}}
\begin{itemize}
  \item \fixed\ \sev{Error} The Prieto et al.\ paper is listed as
    \claim{Correlations in the Data Lead to Semantically Rich Feature
    Geometry under Superposition (ICLR 2026)}. The actual title is
    \emph{From Data Statistics to Feature Geometry: How Correlations Shape
    Superposition}~\cite{prieto2026correlations}. The venue (ICLR 2026) is
    correct. The new bib entry uses the right title.
  \item \fixed\ \sev{Error} Modell et al.: \claim{With the question [\dots] they
    introduce the superposition hypothesis}. They do not introduce it; they
    refer to it as an existing, \enquote{generally agreed upon explanation}
    and cite \textcite{elhage2022toy}. Their own answer to the question is
    different: since layers are linked by linear operations, a curve through
    a higher-dimensional subspace lets linear maps compute more (e.g.,
    polynomials of the feature value)~\cite{modell2025origins}.
  \item \fixed\ \sev{Minor} \claim{Olaf and Batson (2024)} should be \emph{Olah and
    Batson}. Modell et al.\ give the year 2024, although the URL of the
    \emph{Feature Manifold Toy Model} update says 2023, so check before
    citing it.
  \item \fixed\ \sev{Minor} The fourth link for modular arithmetic is
    written as \texttt{(arXiv:2210.01123)} and does not resolve; it should be
    \url{https://arxiv.org/abs/2210.01123}.
  \item \fixed\ \textit{Not an error of the note:} the colour mapping \claim{0: red,
    \dots, $\pi/3$: blue, \dots, $2\pi/3$: yellow} is copied correctly from
    the paper~\cite{modell2025origins}. It does not match the usual hue
    wheel (yellow at $\pi/3$, blue at $4\pi/3$), so do not reuse it without
    the paper's own definition.
\end{itemize}

\paragraph{\note{Methods - finding and using geometry.md}}
\begin{itemize}
  \item \fixed\ \sev{Minor} The GitHub link for Bhalla et al.\ is broken:
    \nolinkurl{https://arxiv.org/pdf/~https://github.com/...}.
    Correct: \url{https://github.com/goodfire-ai/sae-manifold}.
  \item \fixed\ \sev{Inaccurate} \claim{This paper argues that low reconstruction
    error is a trap} reads like a quote, but the paper says \enquote{low
    reconstruction error alone does not guarantee coherent manifold
    recovery}~\cite{bhalla2026saemanifolds}. Use the paper's wording.
  \item \fixed\ \sev{Inaccurate} \claim{strong evidence, that Concepts aren't stored
    in single, isolated features. Instead, a concept is represented by a
    subspace}. The paper makes a narrower claim: for manifolds tiled by an
    SAE, only the \emph{group} of features carries the geometry, so
    interpretation has to reason about subspaces~\cite{bhalla2026saemanifolds}.
    It does not claim this for concepts in general.
\end{itemize}

\paragraph{\note{Benchmarks.md}}
\begin{itemize}
  \item \fixed\ \sev{Inaccurate} MAxBench methods: MFA is listed under
    \claim{linear}. MFA is a mixture of local Gaussian (affine) regions, so
    it is piecewise, not globally linear~\cite{shafran2026regions}. The list
    also misses two of the ten methods, Linear Probe and Factor
    Analysis~\cite{appapogu2026maxbench}.
  \item \fixed\ \textit{Addition:} the limitation of the manifold method is
    stronger than \enquote{ordered concepts}: in its standard form it also
    needs the complete sequence of values, and in the main evaluation it
    is only applied to \emph{Days}~\cite{appapogu2026maxbench}.
\end{itemize}

\paragraph{\note{Foundations - what neural geometry is.md}}
No errors found. The quotes from the Goodfire post and from Engels et al.\
match the sources~\cite{goodfire2026world,engels2025notall}.

\paragraph{\note{Structure.md}}
The writing task is correct: the Goodfire post spells the name
\enquote{Karkada} in footnotes 3 and 11 but \enquote{Kakada} in footnotes 20
and 21~\cite{goodfire2026world}. The exposé cites it as
\textcite{karkada2026symmetry}.

\subsection{Other notes}

\paragraph{\note{unsupervised manifold steering.md}}
\begin{itemize}
  \item \fixed\ \sev{Inaccurate} \claim{Bhalla et al.\ [\dots] read shapes visually
    and don't steer along new ones}. The second half is right. The first
    half is not: on the synthetic benchmark, decoding through the recovered
    atom groups reconstructs topology and geometry \enquote{without
    supervision}~\cite{bhalla2026saemanifolds}. They do steer, but only
    along supervised manifolds (between labelled centroids).
  \item \fixed\ \textit{Update to \claim{Run a fresh literature check}:} the closest
    new work is SMIXAE~\cite{francel2026smixae} (unsupervised discovery with
    a mixture of autoencoders, no steering yet). A search found no paper that
    combines unsupervised discovery with steering along the discovered
    manifold, so the gap still holds (as of \today).
\end{itemize}

\paragraph{\note{Mechanistic Interpretability and Geometry.md} (AI-generated reading list)}
\begin{itemize}
  \item \fixed\ \sev{Error} Elhage et al.\ is called \claim{the mathematical proof of
    how networks exploit high-dimensional geometry}. It is an empirical study
    of small toy models (with some analysis), not a proof about real
    networks~\cite{elhage2022toy}.
  \item \fixed\ \sev{Error} Huh et al.\ (\emph{The Platonic Representation
    Hypothesis}, 2024) \claim{proves mathematically that [\dots] their
    internal vector geometries map out the exact same \enquote{Platonic}
    shape}. It is a position paper: it states a \emph{hypothesis} and gives
    empirical evidence that representations become more similar as models
    scale. It proves no exact identity.
  \item \fixed\ \sev{Error} The Park et al.\ paper is cited as \claim{A Geometric
    Framework for Categorical and Hierarchical Concepts in LLMs}. The actual
    title is \emph{The Geometry of Categorical and Hierarchical Concepts in
    Large Language Models} (Park, Choe, Jiang, Veitch; ICLR 2025).
  \item \fixed\ \sev{Minor} Engels et al.\ is cited as \claim{Not all language model
    features are linear}. The full title is \emph{Not All Language Model
    Features Are One-Dimensionally Linear}~\cite{engels2025notall}.
    \claim{perfectly circular} overstates it: the circles are approximate.
\end{itemize}

\paragraph{\note{Mechanistic Interpretability for AI Safety.md}}
\begin{itemize}
  \item \fixed\ \sev{Error} \claim{Missing in the paper: Mechanistic
    Interpretability (Reverse-Engineering Circuits), Developmental
    Interpretability}. Both are covered by \textcite{bereska2024review}, and
    the same note later summarises them (\enquote{Developmental
    interpretability}, \enquote{Dissecting Models into Interpretable
    Circuits}).
  \item \fixed\ \sev{Inaccurate} \claim{are there alternatives to RL(HF) [\dots]
    yes, [\dots] Representation Engineering}. See the next item.
\end{itemize}

\paragraph{\note{Representation Engineering.md}}
\begin{itemize}
  \item \fixed\ \sev{Inaccurate} The note frames RepE as a replacement for RLHF.
    \textcite{zou2023representation} present it as \enquote{a top-down
    approach to AI transparency}. Steering at inference time is one use, but
    the paper does not position it as an alternative alignment method to
    RLHF. (The old exposé had the same TODO.)
\end{itemize}

\paragraph{\note{Integrated Gradients.md}}
\begin{itemize}
  \item \fixed\ \sev{Error} \claim{Sundararajan et al.\ proved that Integrated
    Gradients is the only attribution method that simultaneously satisfies
    [\dots] Completeness [and] Sensitivity}. The two axioms the method is
    built on are \emph{Sensitivity} and \emph{Implementation Invariance}.
    Completeness is a property it has, and other methods (e.g., DeepLIFT,
    LRP) also satisfy it. The uniqueness result is narrower: IG is the
    unique \emph{path} method that is symmetry-preserving (Sundararajan,
    Taly and Yan, ICML 2017).
  \item \fixed\ \sev{Inaccurate} \claim{solving the saturation problem entirely}.
    IG avoids saturation at the input point, but its result depends on the
    choice of baseline, which is its own well-known problem.
\end{itemize}

\paragraph{\note{Benchmarks for Mechanistic Interpretability.md}}
\begin{itemize}
  \item \fixed\ \sev{Error} SAEBench: \claim{Rather than measuring abstract loss
    metrics (like reconstruction error) [\dots] circuit reconstruction}.
    SAEBench \emph{does} include reconstruction: its first metric,
    \enquote{Core}, covers L0 sparsity and loss recovered. Circuit
    reconstruction is not one of its metrics. The eight metrics are: Core,
    Automated Interpretability, Sparse Probing, RAVEL, Feature Absorption,
    Unlearning, Spurious Correlation Removal and Targeted Probe
    Perturbation~\cite{karvonen2025saebench}.
\end{itemize}

% ===========================================================================
\section{Ratings of the Terminology Notes}\label{sec:ratings}
Rating: should the topic appear in the exposé on unsupervised manifold
steering? 5 = needs its own subsection, 1 = leave it out.

\begin{center}
\small
\begin{tabularx}{\textwidth}{@{}l c X@{}}
\toprule
Note & Rating & Reason \\
\midrule
Topological Data Analysis & 4/5 &
  The note is only a link, but persistent homology is a step of the
  pipeline (finding the shape of a candidate group). It needs a short
  paragraph, e.g., from the old exposé~\cite{gardinazzi2024persistent,tauzin2021giotto}.
  Fill the note before writing it. \\
Representation Engineering & 3/5 &
  One or two sentences as background: linear steering vectors are the
  baseline that manifold steering is compared against~\cite{zou2023representation,wurgaft2026manifoldsteering}.
  The RLHF framing is now corrected (Section~\ref{sec:errors}). \\
Toy Models & 3/5 &
  Superposition is already cited in the Motivation~\cite{elhage2022toy}; a
  citation is enough, no own section. The note now has five sentences on
  \textcite{elhage2022toy} and \textcite{prieto2026correlations}. \\
may trigger normative shifts\dots & 2/5 &
  Hubinger's four ways interpretability helps safety are general
  motivation; one sentence in the Introduction at most. The note's title
  is a sentence fragment from the review; consider renaming it. \\
Benchmarks for Mech.\ Interp. & 2/5 &
  Superseded by \note{Benchmarks.md}; InterpBench and MIB target
  circuits, not manifolds. The SAEBench error is corrected; the note is
  now in \note{unimportant or old}. \\
Integrated Gradients & 1/5 &
  Input attribution; not part of the pipeline. One loose link: IG also
  integrates along a \emph{straight} path, the same assumption manifold
  steering argues against. A footnote at most. \\
SmoothGrad & 1/5 &
  Not relevant to the topic; agree with the \enquote{probably delete} tag. \\
Data & 1/5 &
  Only a link to the MIB datasets, which are for circuit localization. \\
PROBABLY DELETE Memory and Context & 0/5 &
  A different topic (memory compaction, KV caches). Agree to delete. \\
\bottomrule
\end{tabularx}
\end{center}

% ===========================================================================
\section{Bibliography Remarks}\label{sec:bib}
\begin{itemize}[before=\raggedright]
  \item \textit{Completed:} \texttt{shafran2026regions} and
    \texttt{modell2025origins} had only first-author names and short
    titles; both now have full authors and titles.
  \item \textit{Added:} \texttt{goodfire2026world},
    \texttt{karkada2026symmetry}, \texttt{prieto2026correlations},
    \texttt{oozeer2026riemannian}, \texttt{francel2026smixae},
    \texttt{biderman2023pythia}, \texttt{gu2023mamba},
    \texttt{su2021roformer}, \texttt{tenenbaum2000isomap},
    \texttt{coifman2006diffusion}.
  \item \textit{Corrected:} \texttt{korznikov2026sanity} said ICML 2026
    (main conference). It is a poster at the ICML 2026 Workshop on
    Mechanistic Interpretability (\url{https://icml.cc/virtual/2026/79403});
    venue, full authors and full title (\enquote{\dots: Do SAEs Beat Random
    Baselines?}) are now fixed.
  \item \textit{Confirmed:} \texttt{wurgaft2026manifoldsteering} is a
    spotlight at the same workshop (\url{https://icml.cc/virtual/2026/79460}).
  \item \textit{Corrected:} \texttt{gurnee2025manifolds} now cites the
    original Transformer Circuits publication (21~October 2025) with full
    authors and title, and notes the arXiv version (2601.04480, 2026).
  \item \textit{Removed from one claim:} \texttt{bigelow2026stories} was
    cited for the valence--arousal circle, but its abstract is about belief
    updates in story understanding and does not mention emotions. The
    exposé now cites only \texttt{sun2026valence} there.
\end{itemize}

% ===========================================================================
\section{Answers to Next Meeting.md}\label{sec:meeting}
\begin{itemize}
  \item \claim{check if Goodfire used SAEs}: Yes. Bhalla et al.\ (Goodfire)
    train five SAE types (standard, JumpReLU, TopK, BatchTopK, Matryoshka)
    on Llama-3.1-8B, layer 19~\cite{bhalla2026saemanifolds}. Manifold
    steering itself does not use SAEs: it fits a curve through labelled
    activation centroids~\cite{wurgaft2026manifoldsteering}.
  \item \claim{SAEBench: what metrics to use}: all eight score individual
    features or the SAE as a whole under the linear view; none measures
    whether a \emph{group} of features forms a manifold. For this topic,
    Core and RAVEL are the most useful as sanity checks; the
    manifold-specific evaluation has to come from
    \textcite{bhalla2026saemanifolds} and MAxBench~\cite{karvonen2025saebench,appapogu2026maxbench}.
  \item \claim{MAxBench: how do they find/measure the geometries}: ten
    localization methods (DiffMean-r1, ReFT-r1, DiffMean, Linear Probe,
    Schatten Probe, PCA, Factor Analysis, SAE, MFA, spline fitting) recover
    the representation; the benchmark then samples points from it, steers
    with them and scores the generated text for exclusivity (concept
    score) and completeness (diversity score)~\cite{appapogu2026maxbench}.
\end{itemize}

% ===========================================================================
\section{Claude's Ideas Outside the Exposé}\label{sec:ideas}
\begin{itemize}
  \item \textbf{Two-tier benchmark (\note{Ideas.md}) as a by-product.} The
    synthetic tier already exists as the manifold zoo of
    \textcite{bhalla2026saemanifolds}. Instead of a new benchmark, the
    thesis could evaluate its discovery stage on that zoo and its steering
    stage on MAxBench, then report recall on a small catalogue of known
    manifolds in Llama (the real-model tier).
  \item \textbf{Intrinsic interpretability (\note{Ideas.md}).} The idea of
    forcing activations into geometric shapes during training matches the
    \enquote{intrinsic interpretability} branch of
    \textcite{bereska2024review}. It is a separate thesis topic and does
    not fit the current one; keep it for later.
\end{itemize}

\printbibliography[heading=bibintoc,title={References}]

\end{document}
